\section*{The Propers: Detailed Commentary}

\begin{massbox}{Grace gives both order and endurance\enspace\properrefs{Coll., Gosp.}}
The Collect's principal actor is God: he prepares invisible goods, infuses love, and grants that the worshippers may obtain promises surpassing desire. Reconciliation is not a technique by which autonomous wills manufacture communion. It is a fruit of charity received, ordered, and brought into action.
\end{massbox}

\sectionguard
\subsection*{Augustine: The Order of Love\enspace\properrefs{Coll., Gosp., Sec.}}

In \emph{De doctrina christiana} I.22--29 Augustine asks what it means to love God, oneself, the neighbor, and the body in a just order. His vocabulary of ``enjoyment'' and ``use'' can mislead a modern reader unless its destination is retained: God alone is the unchangeable good in whom beatitude rests, while persons are loved in their common relation to him and are helped toward sharing that same good.

\begin{fourstable}{Locus}{Augustine's argument}{Relation to the Collect}{Precise sense}
I.22.20--21 & Neither self nor neighbor is to become the final resting place of desire. The whole current of affection is directed toward God, and the neighbor is loved for God's sake. & ``Above all things'' identifies the unchangeable good; ``in all things'' carries every rightful love into that one current. & Loving a person for God's sake does not license exploitation. The neighbor is truly loved as a person capable of sharing beatitude in God; Augustine's \latin{uti} refers love toward that final good.\\
I.23--26 & Four objects belong within love: God above us, ourselves, the neighbor alongside us, and the body beneath the rational soul. Sin disorders these relations by seeking domination over equals or by mistaking corrupt habit for bodily nature. & Infused charity heals a hierarchy of affections rather than suppressing creaturely goods. & Ordered love is not contempt for the body, personality, or created bonds; it opposes the habit that makes a lower good rule the whole.\\
I.27.28 & A just person loves what ought to be loved and neither gives greater, lesser, nor equal weight contrary to the object's true order. Every human being is loved as human; no sinner is loved precisely as sinner. & The Collect asks for more than intensity. It asks for a truthful measure and direction of love. & Hatred of evil is not hatred of the person. Correcting disorder must remain ordered toward the person's good in God.\\
I.28--29 & All are to be loved, although finite agents cannot give the same concrete help to everyone. Circumstance identifies particular responsibilities; the desired end is that all join in loving God. & The heart can be universal in charity while action remains particular and responsible. & ``All'' does not erase the claims of the actual brother, nor does proximity make that brother the final end of desire.\\
\end{fourstable}

Applied to this Mass, Augustine's order of love bridges the Collect, Gospel, and Secret. The person at the altar cannot appeal to universal benevolence while refusing the particular brother whose grievance has been remembered. Yet repair is not confined to a private dyad: rightly ordered love desires that both persons, and finally all, share the one good who is God.

\begin{studybox}{Neither absorption nor rivalry}
Love of God and neighbor are not two quantities competing for a fixed emotional supply. Augustine says that the current of love for self and neighbor is borne into love of God. The neighbor is not dissolved into an idea of God; he is loved more truly as one who can share in God, and his concrete need lays claim upon the Christian's finite service.
\end{studybox}

\sectionguard
\subsection*{Psalm 26: Interior Cry, Face, and One House\enspace\properrefs{Int., Comm.}}

The Introit and Communion come from the same psalm. Augustine expounds their appointed verses directly, though not their later placement at entrance and Communion.

\begin{fourstable}{Verse and proper}{Augustine's direct exegesis}{Spiritual movement}{Relation to this Mass}
Psalm 26:7, Introit & In \emph{Exposition on Psalm 26}, 7, the voice God hears is interior: a strong intention addressed to his ears and dependent upon mercy. & The worshipper turns affliction into prayer rather than an imagined prosecution of the brother. & Linking the verse to this Mass's coming Gospel is a reading; Augustine does not make it.\\
Psalm 26:8--9, Introit context & The heart seeks not a reward apart from God but God's own countenance; perseverance requires divine help, because the mortal seeks the Eternal as one needing healing. & Desire is purified from using God as an ally for a lesser victory. The face itself becomes the sought good. & Verse 8 stands between the Introit's appointed verses 7 and 9, but its wording is not sung.\\
Psalm 26:4, Communion & The one petition is to remain within the worldwide unity and truth of faith, to behold God's delight face to face, and finally to be made his immortal temple. & The one desired house holds ecclesial perseverance and final vision together. & Augustine directly explains the Communion verse; relating it to reception at this Mass is a reading.\\
\end{fourstable}

The psalm's arc is thus not solitude against community. The heart seeks God himself and, for that very reason, asks to remain within the one house of faith. The Communion antiphon does not allow reconciliation to end as a negotiated ceasefire. Its horizon is shared dwelling, contemplation, and incorporation as God's temple.

\sectionguard
\subsection*{Peter: Blessing and the Pursuit of Peace\enspace\properrefs{Ep.}}

First Peter gives ordered love an ecclesial and bodily form: unity of mind, compassion, brotherly charity, mercy, humility, refusal of retaliation, guarded speech, active goodness, and pursued peace. The apostle then commands the faithful to sanctify Christ in their hearts. Interior allegiance and social action belong to one movement.

\begin{threestable}{Petrine command}{Disorder refused}{Positive form of charity}
Be of one mind, compassionate, loving as brethren & Isolation, indifference, and the assumption that another's suffering is irrelevant. & Ecclesial concord does not require sameness of temperament; it requires shared orientation and suffering-with.\\
Do not return evil or reviling; bless & Retaliation lets the offender determine the Christian's speech and action. & Blessing invokes the good from God and remembers the inherited vocation to blessing.\\
Restrain tongue and lips & Peace cannot coexist with speech designed to deceive, humiliate, or prolong enmity. & Truthful restraint makes room for action directed toward good rather than verbal mastery.\\
Turn from evil, do good; seek peace and pursue it & Peace treated as a mood, a passive absence of conflict, or another person's duty. & The repeated verbs give peace a strenuous shape: departure, good action, search, and pursuit.\\
Sanctify Christ in the heart & A merely horizontal settlement detached from worship. & Christ becomes the heart's governing holiness, from which a meek and truthful answer can proceed.\\
\end{threestable}

\begin{sourcecard}{The Venerable Bede --- on 1 Peter 3}
In \emph{Super epistolas Catholicas}, \emph{In Primam Epistolam Petri} III, Bede reads Peter's differentiated Christians as gathered into one heart in the Lord's faith. Blessing already rehearses their inheritance; persecutors may inflict loss or torture but cannot take the final good from those who persevere; sanctifying Christ inwardly gives courage to explain hope with humility and a good conscience. Bede's lemma reads \latin{In fide autem}, not the Mass text's \latin{in oratione estote}: he comments directly on the chapter in a Latin form that differs from the Mass text there.
\end{sourcecard}

Peter's words ``seek peace and pursue it'' quote Psalm 33. Jerome's \emph{Letter 125}.14 therefore receives the same psalm directly and does not comment on First Peter. Jerome insists that leaving evil requires doing its contrary good, and that peace must be chased even when it seems to flee. His monastic setting differs from this Sunday, but the active force of ``pursue'' is exact.

Anthony of Padua's preceding-Sunday sermon uses this very Petrine lesson beside Luke's catch. At 1~Peter 3:13--15 he distinguishes suffering for justice from suffering for fault, refuses fear that destabilizes the interior person, and makes sanctifying Christ in the heart the condition for truthful witness. Bede follows the letter's syntax more continuously; Anthony shows how a medieval preacher joined courage, imitation of holy teachers, and Peter's own humbled call without turning persecution into proof of innocence.

\begin{studybox}{Peace is not moral vagueness}
Peter's sequence joins peace to truthful speech, departure from evil, positive good, and a clean conscience. The command does not call falsehood true, erase justice, or make coerced contact the measure of forgiveness. It does forbid cherished retaliation and indifference to repair. Prudence may govern the manner of approach; it cannot baptize enmity as peace.
\end{studybox}

\sectionguard
\subsection*{The Gradual: A Psalm Source Recast for the Servants\enspace\properrefs{Grad.}}

The printed reference is Psalm 83:10 and 9, but the Gradual is not a verbatim biblical excerpt. It reverses the source order and changes the object of each petition. The resulting plural language is therefore proper to the liturgical chant: the Korahite psalm does not contain it, and it is not the wording Augustine expounded.

\begin{threestable}{Unit}{Vulgate source}{Missal form and effect}
Gradual response, from verse 10 & \latin{Protector noster, aspice, Deus, et respice in faciem christi tui.} & The opening is retained, but \latin{in faciem christi tui} becomes \latin{super servos tuos}: God is asked to look upon his servants rather than upon the face of his anointed.\\
Versicle, from verse 9 & \latin{Domine Deus virtutum, exaudi orationem meam; auribus percipe, Deus Jacob.} & The address and \latin{exaudi} remain, but the singular prayer and appeal to Jacob's God become \latin{preces servorum tuorum}, again a plural servant petition.\\
\end{threestable}

The complete Psalm is a pilgrimage prayer moving from longing for God's courts through ascent toward Sion, then to hearing, the shield or anointed king, and confidence in the Lord. Augustine, \emph{Exposition on Psalm 84}, 12--13, follows that unadapted sequence: he asks that Jacob be made Israel and reads the face of God's Christ as Christ made known. His exposition directly receives the source verses and does not comment on the Gradual's substitutions. After the Epistle, the adapted chant gives the assembled servants a common request to be seen and heard.

\sectionguard
\subsection*{The Alleluia: Royal Joy Is Received\enspace\properrefs{All.}}

Psalm 20 is a royal thanksgiving. Its opening gives the king's joy two divine objects---the Lord's strength and salvation---before the complete psalm proceeds through granted desire, crown, life, trust, victory, and judgment. The Alleluia selects the opening acclaim, so its joy is neither self-congratulation nor evidence that royal power validates itself.

\begin{threestable}{Reading level}{Contribution}{Precise sense}
Psalm in context & Israel's anointed king gives thanks because strength and salvation come from the Lord; the remaining verses place royal desire and endurance under divine gift. & The particular king, occasion, and composition history are not recoverable from the verse.\\
Augustine, \emph{Exposition on Psalm 21}, 1--2 & Augustine directly reads the king as Christ, the Lord's strength through the Word made flesh, and salvation as God's life-giving work. & Augustine's reading is allegorical, and the royal psalm keeps its literal and historical horizon.\\
Alleluia before the Gospel & The assembly acclaims received strength and salvation immediately before hearing the King's command concerning righteousness and reconciliation. & The liturgical adjacency supports attention to Christ's authority; it does not turn every detail of the psalm into commentary on Matthew 5.\\
\end{threestable}

\sectionguard
\subsection*{Matthew 5: From the Root to the Altar\enspace\properrefs{Gosp.}}

Matthew 5:20--24 belongs to Christ's account of the righteousness that exceeds that of the scribes and Pharisees. The movement is inward and then sacramental: from homicide to anger, from anger to contemptuous address, from speech to remembered grievance, and from remembered grievance to interrupted offering and urgent reconciliation.

\begin{fourstable}{Witness and locus}{Direct reading of the Gospel}{Sacrificial claim}{Interpretive precision}
St. John Chrysostom, \emph{Homily 16 on Matthew}, 6, 9--12 & The greater righteousness does not reject the old law but brings it to stricter perfection. Contemptuous speech matters because small words feed the fire of enmity and can grow into greater violence. & Christ permits divine service to be interrupted because reconciled charity is itself sacrifice; without charity the other gift is not received. & Chrysostom first addresses the wrongdoer, then widens the urgency so that the injured person also refuses prolonged enmity. This does not make guilt symmetrical.\\
St. Augustine, \emph{On the Sermon on the Mount} I.9.21--I.10.28 & The linked sayings trace anger retained in the mind, contempt expressed, and the humble remedy that cuts off hatred. ``Your brother has something against you'' directs remembrance toward injury one has caused. & Augustine reads the visible altar together with faith in God's interior temple: the offerer must approach with a mind free from dissimulation and ask pardon where injury was done. & His discussion of an absent brother prevents a wooden demand to abandon a gift indefinitely; interior humility must be real and should issue in truthful approach when possible.\\
\end{fourstable}

Chrysostom protects the objective seriousness of the altar: Christ does not abolish sacrifice but refuses its use as refuge from a ruptured charity. Augustine protects the direction of examination: the text first makes the offerer remember what the brother has \emph{against him}, not construct a catalogue of the brother's offenses. Together they turn memory from accusation to responsibility.

Anthony's sermon for the Sixth Sunday after Pentecost reads the passage with an explicitly penitential medieval emphasis. Greater righteousness reaches the intention and not only the hand; anger is a seed whose consent can mature into deadly passion; and the gift at the altar cannot substitute for a spirit still stained by injury. His moralized ``brother'' sometimes becomes the offerer's own spirit, while the Gospel's command concerns a literal brother. His examination is nevertheless sharp: bodily austerity and external worship cannot purchase exemption from inward truth or actual reconciliation.

\begin{massbox}{The altar reveals rather than conceals}
The gift lying before the altar makes delay visible. The worshipper cannot purchase innocence by offering, but neither is the altar opposed to reconciliation. Christ orders the offering by ordering the offerer toward the brother, so that sacrificial worship and charity may become one truthful act.
\end{massbox}

\sectionguard
\subsection*{The Secret: What Each Offers Profits All\enspace\properrefs{Sec.}}

The Secret does not dissolve individual agency into an anonymous collective. It names individual offering distributively and then asks for universal ecclesial fruit.

\begin{threestable}{Latin hinge}{Grammatical force}{Ecclesial movement}
\latin{has oblationes famulorum famularumque tuarum benignus assume} & God is asked mercifully to receive these oblations of his servants, women and men. & Acceptance remains divine gift; offering does not validate itself.\\
\latin{quod singuli obtulerunt} & The neuter singular \latin{quod} gathers the offered reality, while distributive \latin{singuli} preserves ``each one.'' & Personal histories, responsibilities, and acts are not erased at the altar.\\
\latin{ad honorem nominis tui} & The honor of God's Name is the offering's governing reference. & Augustine's order of love reaches sacrificial form: the gift is not directed to private prestige or control.\\
\latin{cunctis proficiat ad salutem} & The prayer asks that it profit everyone unto salvation. & What each gives is received into an end no isolated offerer can possess privately.\\
\end{threestable}

\begin{sourcecard}{St. Cyprian --- \emph{On the Lord's Prayer}, 8}
Cyprian does not comment on this Secret, but his teaching on common prayer bears closely on it. Christian prayer says ``our,'' ``us,'' and ``we'' because the people are one; Christ teaches each Christian not to pray for self alone but for all. Individual voices are not silenced. Like the Three Children praying as from one mouth, they become a concordant common petition. This closely clarifies the Secret's passage from \latin{singuli} to \latin{cunctis}.
\end{sourcecard}

\begin{sourcecard}{St. Augustine --- \emph{City of God} X.5--6}
Augustine teaches sacrificial doctrine and does not expound the oration. True sacrifices are works of mercy ordered to holy fellowship with God; body and soul are offered as the mind is transformed; and the whole redeemed city is offered through Christ the High Priest. The Secret's many individual offerings can therefore serve all because the offerers are gathered into an ecclesial sacrifice whose unity is Christ's.
\end{sourcecard}

\sectionguard
\subsection*{The Offertory: Counsel, Presence, and Steadiness\enspace\properrefs{Off.}}

Psalm 15 begins with refuge, chosen inheritance, and refusal of rival gods. Its appointed verses then bless the Lord for counsel, keep him continually before the psalmist, and locate him at the right hand so that the speaker will not be moved. The verses immediately following confess bodily security, deliverance from corruption, the path of life, and joy in God's presence. The Offertory sings only verses 7--8: its explicit action is grateful attention and God-given steadiness, not a declaration that every decision of the offerer is already just.

Acts 2:25--32 later quotes Psalm 15:8--11 in Peter's proclamation of Christ's resurrection. That canonical reception gives the chant a genuine paschal horizon: the Lord kept before the eyes and at the right hand is joined to the Holy One not abandoned to corruption. Yet the Mass does not print the later verses or Peter's interpretation in this chant. At the offering, the later verses and Acts therefore supply context and a canonical horizon, and the appointed words remain verses 7--8.

Augustine's direct exposition of verses 7--8 speaks in Christ's assumed humanity: understanding is received from the Lord, and the passing things of mortal history never remove the speaker's sight from the One who abides. In Augustine's \emph{totus Christus} horizon, the Church shares that fixed attention through her Head. The reading supplies more than generic confidence, but its Christological voice is Christ's and does not claim Christ's sinlessness for each offerer.

\begin{threestable}{Appointed action}{Larger biblical horizon}{Precise sense}
The offerer blesses the Lord who gives counsel. & Psalm 15 places every other inheritance beneath the Lord, the psalmist's portion. & Divine counsel does not certify the offerer's private account of a conflict.\\
The Lord is kept always in sight and at the right hand. & The psalm proceeds from steadiness to secure rest, the path of life, and joy in God's presence. & The later verses illuminate the selected couplet but are not themselves sung.\\
Acts applies verses 8--11 to Christ risen from the dead. & The Offertory's steadiness can be heard within a received paschal reading of its complete psalm. & The apostolic proof concerns Christ; it does not turn the offerer's confidence into personal impeccability.\\
\end{threestable}

\sectionguard
\subsection*{Communion and Postcommunion: Fullness Still Asks\enspace\properrefs{Comm., Postcomm.}}

Augustine reads the Communion's one desire and one house as the worldwide unity of faith and the future temple. The Postcommunion now answers that concentrated desire from within accomplished reception. Its opening relative clause, \latin{Quos caelesti, Domine, dono satiasti}, addresses the Lord and identifies the people as those whom \emph{he} has satiated with a heavenly gift. The governing petition, \latin{praesta, quaesumus}, then asks God to grant what follows. The gift is received fullness; it is not the sentence's grammatical agent.

\begin{threestable}{Clause}{What the prayer says}{Doctrinal sense}
\latin{Quos caelesti, Domine, dono satiasti} & God has already filled the addressed community with a heavenly gift. & Reception is acknowledged as an accomplished divine action, not reduced to aspiration or feeling.\\
\latin{ut a nostris mundemur occultis} & The recipients ask to be cleansed from faults hidden within their own life. & Being filled does not imply complete self-knowledge or finished conversion.\\
\latin{et ab hostium liberemur insidiis} & They also ask to be delivered from hostile snares outside them. & Interior fault and exterior hostility are coordinated without confusing guilt, danger, or responsibility.\\
\end{threestable}

The paired passive subjunctives \latin{mundemur} and \latin{liberemur} keep communicants recipients of God's action.

\Needspace{3\baselineskip}
\noindent The checked patristic witnesses supply no secure direct commentary on this oration. In its complete Latin grammar and its place after Communion, the sequence from one house to further cleansing does not diminish the gift: it denies only that sacramental fullness makes vigilance, conversion, or dependence obsolete.

\sectionguard
\section*{The Propers: Notable and Quotable}

{\small
\begin{itemize}[itemsep=0.35em,topsep=0.2em,parsep=0pt]
\item \textbf{Introit, Psalm~26:1 --- University of Oxford arms and Vice-Chancellor's Oration (2023).} The unaccented \latin{Dominus illuminatio mea} has served as Oxford's legend since at least the later sixteenth century. Irene Tracey's oration then dropped a few letters to make \latin{Domus illuminat nos}---``our home enlightens us.'' The psalmist's confession thus becomes both a corporate badge and a self-aware collegiate pun, shifting the grammatical light from the Lord to the university community without pretending to alter the crest.

\item \textbf{Gospel, Matthew~5:22 --- Honoré de Balzac, \emph{The Physiology of Marriage}, Meditation II, ``Marriage Statistics'' (1829).} In a satirical pseudo-census, leisured aesthetes exclude laboring people from their idea of fully human life: ``We say Raca and fool'' to anyone not young, beautiful, and passionate. Christ's condemned insults become the privileged speakers' own vocabulary; the allusion makes their cultivated contempt indict itself rather than endorsing it.

\item \textbf{Offertory, Psalm~15:8 --- Shimon Zimel, \emph{Shiviti} (1845/46), Jewish Museum accession JM~51-51.} ``I have set the Lord always before me'' moves from first-person resolve into a synagogue furnishing placed before the Torah reader. Zimel's cut-parchment plaque surrounds the verse with dense scrollwork, vegetation, and real and mythical animals: inward recollection becomes a literal field of vision, not merely a sentence quoted in prose.
\end{itemize}
}

\sectionguard
\section*{The Propers: Interpretive Possibilities}

\noindent\emph{The following exploratory proposals were developed in an AI-assisted editorial process. They are attributed to none of the cited authorities and claim no historical compositional intent. Each states its anchors and limits.}

\Needspace{8\baselineskip}
\noindent\textbf{The sought face becomes face-to-face responsibility.\enspace\properrefs{Int., Coll., Gosp., Comm.}}
The Introit seeks God's face and asks to be heard; the Collect asks that every love be ordered; the Gospel interrupts altarward movement with a remembered brother; Communion gathers desire into dwelling in God's house. The mechanism is a devotional trajectory in which vertical attention is tested by a lateral responsibility before returning to common worship. Read separately, seeking God and examining a human grievance can remain parallel duties; the sequence proposes that contemplation is falsified when it becomes a route around the concrete brother whose claim has been recalled. Its spiritual fruit is an examen of whether prayer enlarges truthful presence to God and neighbor or merely protects a religious self-image. The grievance's justice and the offerer's culpability remain to be established; the divine face is not identified with the neighbor, and prudence may require investigation, distance, safeguarding, restitution, or mediated contact rather than immediate encounter.

\Needspace{8\baselineskip}
\noindent\textbf{Blessing breaks retaliation's closed circuit.\enspace\properrefs{Ep., Grad., Gosp., Off.}}
First Peter replaces returned evil with blessing; the Gradual transfers the demand to be seen and heard to the protecting God; the Gospel requires initiative from the one who remembers a grievance; the Offertory fixes the Lord before the eyes. The connecting mechanism is a redistribution of speech, attention, and initiative: harm no longer supplies the script for the next word or act. An element-by-element reading can miss how blessing is neither passivity nor verbal niceness but the opening of a closed retaliatory circuit toward divine judgment and reparative action. The pastoral fruit is a sequence for conflict discernment: refuse mimetic injury, seek God's hearing, name one's responsibility, act toward repair, and keep judgment under God's gaze. Its strongest limit is that this neither equalizes culpability nor suppresses protection and justice, and reconciliation cannot be manufactured by one party.

\Needspace{8\baselineskip}
\noindent\textbf{The altar opens a breach in selective memory.\enspace\properrefs{Gosp., Sec., Postcomm.}}
\enlargethispage{\baselineskip}
In the Gospel, arrival at the altar awakens memory of another's claim; the Secret makes each person an offerer; after reception, the Postcommunion asks cleansing from hidden faults. The mechanism is a sacramental sequence from interrupted ritual action, through personal oblation, to confessed limits of self-knowledge. Read separately, the Gospel can become a preliminary rule and the later prayers routine completion; together they suggest that worship does not certify the worshipper's innocence but places memory, gift, and opacity before God. The intellectual and spiritual fruit is a liturgical account of conscience in which recollection can begin before the gift yet purification remains necessary afterward. The disconfirming limit is that ritual does not automatically reveal guilt, every allegation is not true, and hidden faults cannot be named confidently without evidence, honest examination, and due process.

\Needspace{8\baselineskip}
\noindent\textbf{The singular gift acquires a plural horizon.\enspace\properrefs{Coll., Sec., Comm.}}
The Collect orders all loves under God; the Secret asks that what \latin{singuli} offer profit \latin{cunctis}; Communion voices one desire for one divine house. The mechanism is a grammatical and spatial expansion: distinct offerers remain distinct, but God redirects their gifts toward a common salvation and dwelling. A proper-by-proper reading can miss the ecclesial transformation between ``each'' and ``all,'' treating offering as private intention and Communion as private consolation. The fruit is a criterion for charisms, labor, patronage, and institutional resources: a gift matures when it can be received, tested, and made fruitful beyond the giver's possession. The strongest limit is that common benefit does not erase particular duties, intentions, consent, or governance, and a private project does not acquire ecclesial authority merely because its sponsor calls it an offering.

\Needspace{8\baselineskip}
\noindent\textbf{The reconciler remains a patient of reconciliation.\enspace\properrefs{Grad., Gosp., Off., Postcomm.}}
The Gradual's servants ask God to look and hear; the Gospel gives those dependent servants a demanding initiative toward repair; the Offertory steadies their gaze on the Lord; the Postcommunion returns the same plural voice to hidden cleansing and defense. The mechanism is pronoun continuity across petition, agency, attention, and renewed need. Read in isolation, ``go first'' can be heard as the morally competent person managing another's defect; the complete movement instead places the reconciler inside the mercy being enacted. Its pastoral fruit is courageous initiative without superiority: one can make restitution or seek peace precisely as a person still needing correction and protection. The limit is that shared dependence does not equalize wrongs, suspend judgment, require unsafe contact, or excuse the party with the greater responsibility from making the first concrete repair.
